% Commutative diagram from the long-exact-sequence chapter.
% Source: book/tex/homology/long-exact.tex (lines 128-134).
\documentclass[border=2mm]{standalone}
\usepackage{amssymb,amsmath}
\usepackage{tikz-cd}
\newcommand{\ZZ}{\mathbb{Z}}
\newcommand{\eps}{\varepsilon}
\newcommand{\id}{\mathrm{id}}
\begin{document}
\begin{tikzcd}
	\vdots & \vdots & \vdots  \\
	H_{n+1}(A_\bullet) \ar[r, "f_\ast"] & H_{n+1}(B_\bullet) \ar[r, "g_\ast"] & H_{n+1}(C_\bullet) \\
	H_{n}(A_\bullet) \ar[r, "f_\ast"] & H_{n}(B_\bullet) \ar[r, "g_\ast"] & H_{n}(C_\bullet) \\
	H_{n-1}(A_\bullet) \ar[r, "f_\ast"] & H_{n-1}(B_\bullet) \ar[r, "g_\ast"] & H_{n-1}(C_\bullet) \\
	\vdots & \vdots & \vdots \\
\end{tikzcd}
\end{document}
